\scanpage{45}
\mspage{132}

Notons que l'hyp sur $S_0$ assure que si $F,F'$ sont
deux faisceaux d'ens. loc. const. sur $S$ et $u:F\unsure{\text{---}|\to}F'$
un \deleted{\unsure{h}}hom \deleted{qui \unsure{sont|est} iso \unsure{sur}} \deleted{\unsure{un}} tel que $u_s:F_s\simeq F'_s$
pour $s\in\unsure{S|S_0}$, alors $u$ est iso.

$1^\circ$) Soit $F$ un faisceau d'ens \inserted{loc} constant sur $X$,
$F'$ son image inverse sur $X'$, il faut prouver
\[
H^0(X,F)\xrightarrow{\sim}H^{\unsure{0|\bullet}}(X',F').
\]
Or
\[
H^0(X,F)\simeq H^0(Y,\underbrace{\unsure{p|f}_*(F)}_{\substack{\text{loc. const}\\\text{sur }S}})\longrightarrow H^0(Y,\unsure{p'|f'}_*(F'))\simeq H^0(X',F'),
\]
il faut prouver que \inserted{$\unsure{(*)|(x)}\quad p_{\unsure{*|=}}(F)\xrightarrow{\sim}p'_*(F')$}
est iso. Or par hyp. sur $p,p'$, les deux
faisceaux sont loc. constants et se calculent
fibre par fibre, d'où résulte que (*) est un
isom pour les fibres en $s\in S_0$, donc un iso
partout. OK

$2^\circ$) Il en résulte déjà, \inserted{si $F$ est un faisceau en \uline{groupes},} que le foncteur
\[
F\text{-torseurs sur }X\text{---}F\text{-torseurs sur }X'
\]
est pl. fidèle, donc $H^1(X,F)\to H^1(X',F')$
injectif, il \deleted{\unsure{f}} reste à prouver que ce foncteur
est ess$^{\mathrm t}$. surjectif. À cause de la pleine
fidélité, c'est \uline{local sur $S$} (\uline{NB} On voit directement que l'hyp. faite sur les $s\in S_0$ implique
la $=$ hyp. pour \uline{tout} $s\in S$, en utilisant le
fait que $p,p'$ \inserted{sont} des fibrations.). Donc on est ramené
au cas où $S=S_{/s_0}$ a un objet final $s_0$.

% slanted
\margin{Les hyps faites sont bel et bien locales sur $S$ (du moins si $S_0=S$)}

% archive sheet number: 45
